%! Setting Up --- Learning from Data — Unit 8, Lesson 8.0 — Exit Ticket
\documentclass[10pt]{article}
\usepackage{linalg-article}
\usepackage{linalg-boxes}
\newcommand{\TallMath}[1]{$\displaystyle #1\rule[-1.4em]{0pt}{3.2em}$}
\newcommand{\vv}{\mathbf{v}}
\newcommand{\relu}{\mathrm{ReLU}}

\begin{document}

\pageheader{Unit 8, Lesson 8.0}{Exit Ticket}
\namedateperiod

\vspace{0.15in}

No notes. Show your reasoning.

\begin{enumerate}[leftmargin=1.8em, itemsep=16pt]
  \item \textbf{One layer.} For $A=\begin{bmatrix} 1 & -2 \\ 1 & 1 \end{bmatrix}$ and
        $\vv=\begin{bmatrix} 3 \\ 2 \end{bmatrix}$, compute the linear step $A\vv$, then apply the ReLU to
        each entry to get $\relu(A\vv)$.
        \vspace{1.3cm}
  \item \textbf{Measure the error.} A network predicts $2,5,3$ while the true targets are $3,5,1$. Compute
        the loss (the sum of squared errors).
        \vspace{1.2cm}
  \item \textbf{Explain.} In one sentence, why does a neural network need the nonlinear ReLU step ---
        what goes wrong if every step is just a matrix multiplication?
        \writelines{2}
\end{enumerate}

\end{document}
